\documentclass[11pt,a4paper]{article}
\usepackage[utf8]{inputenc}
\usepackage[T1]{fontenc}
\usepackage{amsmath,amsfonts,amssymb}
\usepackage{graphicx}
\usepackage{hyperref}
\usepackage{listings}
\usepackage{color}
\usepackage{booktabs}
\usepackage{float}
\usepackage{geometry}
\geometry{margin=1in}
\definecolor{zenblue}{RGB}{41,121,255}
\hypersetup{colorlinks=true,linkcolor=zenblue,urlcolor=zenblue,citecolor=zenblue}

\title{\textbf{Reducing Hallucinations in Zen Models: Retrieval Augmentation,\\
Confidence Calibration, and Citation-Grounded Generation}\\
\large Technical Report v2025.08}
\author{Antje Worring, Zach Kelling \\ Zen LM Research Team\\
\texttt{research@zenlm.org}}
\date{August 2025}

\begin{document}
\maketitle

\begin{abstract}
Hallucination — the generation of plausible but factually incorrect content — remains
a central challenge for large language models. We report a comprehensive study of
hallucination in the Qwen3-based Zen models (dense 0.6B/4B/8B/32B and the Qwen3-30B-A3B
mixture-of-experts variant, all Apache-2.0) and introduce three complementary
interventions: (1) \textbf{Retrieval-Augmented Factuality Training (RAFT)}, which
fine-tunes models to condition on retrieved evidence when generating factual claims;
(2) \textbf{Uncertainty-Aware Decoding (UAD)}, which modulates token probabilities by
a calibrated confidence estimate to suppress low-confidence continuations; and (3)
\textbf{Citation-Grounded Generation (CGG)}, which trains models to emit inline
citations linking claims to source documents. Together, these methods substantially
reduce hallucination rates on TruthfulQA, FactScore, and HaluEval relative to the
unmodified base models, with negligible perplexity degradation. We report the
\emph{relative} contribution of each component rather than restating third-party
benchmark figures for the base models.
\end{abstract}

\tableofcontents
\newpage

%% -----------------------------------------------------------------------
\section{Introduction}
\label{sec:intro}
%% -----------------------------------------------------------------------

Language models learn to predict the next token by compressing statistical patterns in
training text. When queried about facts outside their training distribution, or about
subtle factual nuances within it, they may generate fluent but incorrect content with
high confidence. This phenomenon — hallucination — erodes user trust and limits
deployment in high-stakes domains such as medicine, law, and scientific research.

Hallucination manifests in several forms:
\begin{itemize}
    \item \textbf{Factual hallucination}: asserting a false proposition about the world
          (e.g., incorrect dates, attributed quotes, fabricated citations).
    \item \textbf{Knowledge boundary hallucination}: confidently answering questions
          that are genuinely unanswerable from training data.
    \item \textbf{Reasoning hallucination}: correct premises leading to incorrect
          conclusions through flawed inference steps.
\end{itemize}

We address all three categories through our RAFT + UAD + CGG system, evaluated on the
Qwen3-based Zen models across the 4B, 8B, and 32B dense scales (and the Qwen3-30B-A3B
MoE variant).

\subsection{Contributions}

\begin{enumerate}
    \item A systematic measurement of hallucination rates across domains and model scales,
          using TruthfulQA, FactScore, HaluEval, and a new domain-stratified hallucination
          benchmark (ZenHallu).
    \item Retrieval-Augmented Factuality Training (RAFT): a fine-tuning recipe that
          teaches models to identify when retrieval is needed and how to synthesize
          evidence faithfully.
    \item Uncertainty-Aware Decoding (UAD): a training-free decoding modification that
          calibrates output confidence and suppresses hedging-unfaithful generations.
    \item Citation-Grounded Generation (CGG): a training procedure for models to emit
          structured inline citations with confidence scores.
    \item Ablation and analysis showing complementary contributions of each component.
\end{enumerate}

%% -----------------------------------------------------------------------
\section{Hallucination Measurement}
\label{sec:measurement}
%% -----------------------------------------------------------------------

\subsection{Benchmark Descriptions}

\textbf{TruthfulQA} \cite{lin2021truthfulqa}: 817 questions spanning conspiracy theories,
misconceptions, fiction, and factual recall. Models must answer truthfully or abstain.
We use the multiple-choice formulation.

\textbf{FactScore} \cite{min2023factscore}: Measures the factual precision of long-form
biography generation by decomposing outputs into atomic claims and verifying each against
Wikipedia. Score $\in [0, 1]$, higher is better.

\textbf{HaluEval} \cite{li2023halueval}: A large-scale hallucination evaluation benchmark
with 35,000 samples across QA, dialogue, and summarization. Each sample contains a
human-annotated hallucinated response and a correct response; models are evaluated on
classification accuracy.

\textbf{ZenHallu} (ours): A domain-stratified benchmark covering medicine, law, science,
history, mathematics, and geography. Each domain contains 500 factual questions with
verified ground truth. Hallucination rate is measured as the fraction of generated
responses containing at least one verifiable factual error.

\subsection{Baseline Hallucination Rates}

Measuring the unmodified Qwen3-based Zen models (4B, 8B, and 32B dense) on the four
benchmarks above, we observe the expected scale trend: larger models hallucinate less,
scoring higher on TruthfulQA and FactScore and lower on the ZenHallu error rate, with the
32B dense model the strongest of the dense variants. These baseline measurements
establish the reference point against which the RAFT + UAD + CGG interventions are
evaluated in Section~\ref{sec:experiments}; we report intervention effects as
improvements relative to this baseline rather than as standalone capability claims.

\subsection{Domain Stratification}

Stratifying ZenHallu by domain (each domain contains 500 verified factual questions),
hallucination rates on the 32B baseline are highest in \textbf{medicine} and
\textbf{law}, intermediate in \textbf{science} and \textbf{history}, and lowest in
\textbf{geography} and \textbf{mathematics}. This ordering reflects the relative
density of verifiable, frequently-updated factual claims in each domain: specialized,
fast-changing knowledge (clinical facts, statutes) is harder to recall reliably than
stable, heavily-represented knowledge (geographic facts, arithmetic).

%% -----------------------------------------------------------------------
\section{Retrieval-Augmented Factuality Training (RAFT)}
\label{sec:raft}
%% -----------------------------------------------------------------------

\subsection{Motivation}

Standard retrieval-augmented generation (RAG) appends retrieved context to the input
prompt at inference time. This helps if the model knows to trust the retrieved context,
but baseline models often ignore retrieved evidence or confabulate details not present
in it. RAFT trains the model to faithfully condition on retrieved evidence.

\subsection{Training Procedure}

RAFT fine-tunes on a dataset of (question, retrieved documents, answer) triples, where
the answer is grounded in the retrieved documents. We construct the RAFT training set
using:

\begin{enumerate}
    \item \textbf{Retrieval}: For each factual question $q$, retrieve the top-$k$ documents
          $\{d_1, \ldots, d_k\}$ from a 100B-token knowledge base using a dense retriever.
    \item \textbf{Distractor injection}: With probability $p_{\text{dist}} = 0.4$, replace
          one retrieved document with a semantically similar but factually incorrect
          distractor. This trains the model to critically evaluate rather than blindly copy.
    \item \textbf{Chain-of-thought grounding}: The target answer includes an explicit
          chain of evidence linking claims to specific document spans, formatted as
          ``[Source: $d_j$, span: ...] $\Rightarrow$ claim''.
\end{enumerate}

\subsection{RAFT Loss}

The RAFT fine-tuning loss combines standard language modeling with a faithfulness penalty:
\begin{equation}
    \mathcal{L}_{\text{RAFT}} = \mathcal{L}_{\text{CE}} + \beta \cdot \mathcal{L}_{\text{faith}}
    \label{eq:raft_loss}
\end{equation}

The faithfulness loss $\mathcal{L}_{\text{faith}}$ penalizes generated claims that cannot
be grounded in the retrieved documents. It is computed using an NLI classifier $h$ that
scores the entailment between each generated claim $c_t$ and the document set:
\begin{equation}
    \mathcal{L}_{\text{faith}} = \frac{1}{T}\sum_{t=1}^{T}
    \max\!\bigl(0,\; \gamma - \max_j h(d_j, c_t)\bigr)
    \label{eq:faith_loss}
\end{equation}

where $\gamma$ is a faithfulness margin and $T$ is the number of generated claims.

%% -----------------------------------------------------------------------
\section{Uncertainty-Aware Decoding (UAD)}
\label{sec:uad}
%% -----------------------------------------------------------------------

\subsection{Confidence Calibration}

A well-calibrated model should express uncertainty proportional to the likelihood that
its output is correct. We measure calibration via the Expected Calibration Error (ECE):
\begin{equation}
    \text{ECE} = \sum_{m=1}^{M} \frac{|B_m|}{n} \left|
    \text{acc}(B_m) - \text{conf}(B_m) \right|
    \label{eq:ece}
\end{equation}

where $B_m$ partitions predictions by confidence into $M$ bins.

The baseline Zen-32B model shows a substantial Expected Calibration Error, indicating
significant overconfidence that UAD is designed to correct.

\subsection{Temperature Scaling with Semantic Adjustment}

UAD applies a two-stage calibration: first, global temperature scaling $T^*$ learned on
a held-out calibration set to minimize ECE; second, a semantic confidence adjustment that
modulates the effective temperature based on the model's uncertainty about the semantic
content of the response:

\begin{equation}
    p_{\text{UAD}}(y_t \mid y_{<t}, x) \propto
    \exp\!\left(\frac{\log p(y_t \mid y_{<t}, x)}{T^* \cdot (1 + \kappa \cdot u_t)}\right)
    \label{eq:uad}
\end{equation}

where $u_t \in [0,1]$ is a per-token uncertainty estimate derived from the model's
entropy over the top-$V$ vocabulary candidates, and $\kappa$ controls the strength of
the adjustment. High $u_t$ increases the effective temperature, producing a flatter
distribution that prevents overconfident hallucinated tokens.

\subsection{Abstention Mechanism}

When the aggregate uncertainty $\bar{u} = \frac{1}{T}\sum_t u_t$ exceeds a threshold
$\theta_{\text{abs}}$, the model emits an abstention token ``I am not confident enough
to answer this question'' rather than a potentially hallucinated response. This is
calibrated to achieve a false abstention rate below 5\% on answerable questions.

%% -----------------------------------------------------------------------
\section{Citation-Grounded Generation (CGG)}
\label{sec:cgg}
%% -----------------------------------------------------------------------

\subsection{Citation Format}

CGG trains models to interleave inline citations within generated text:
\begin{verbatim}
The boiling point of water at sea level is 100°C [cite:doc_42, p.7]
with a boiling point elevation of 0.512°C·kg/mol [cite:doc_18, p.14].
\end{verbatim}

Each citation specifies a document ID and page/paragraph range from the retrieved context.

\subsection{Citation Training}

CGG adds a citation prediction head to the model that, at each position where a factual
claim is completed, predicts the most relevant source span. The head is a lightweight
cross-attention module over the retrieved document embeddings:
\begin{equation}
    \alpha_j = \text{softmax}\!\left(\frac{h_t^\top D_j}{\sqrt{d}}\right), \quad
    \text{cite}_t = \operatorname{argmax}_j \alpha_j
    \label{eq:cite}
\end{equation}
where $h_t$ is the model's hidden state at position $t$ and $D_j$ is the embedding of
retrieved document chunk $j$. The citation head is trained with a cross-entropy loss
against human-annotated citation labels.

\subsection{Citation Confidence Score}

Alongside the cited document, the model emits a confidence score $c_t = \max_j \alpha_j$.
Low citation confidence is correlated with hallucination and can be surfaced to users
as a reliability indicator.

%% -----------------------------------------------------------------------
\section{Experiments}
\label{sec:experiments}
%% -----------------------------------------------------------------------

\subsection{Main Results (Component Ablation)}

\begin{table}[H]
\centering
\caption{Direction of effect of each anti-hallucination component on the Zen-32B
(Qwen3-32B) model. ``$\uparrow$'' indicates improvement on a benchmark (lower hallucination
/ higher factuality); more arrows indicate a larger relative effect. The full
RAFT + UAD + CGG system is strongest on every benchmark.}
\begin{tabular}{lcccc}
\toprule
\textbf{System} & \textbf{TruthfulQA} & \textbf{FactScore} & \textbf{HaluEval} & \textbf{ZenHallu} \\
\midrule
Baseline               & ---  & ---  & ---  & --- \\
+ RAFT                 & $\uparrow\uparrow$   & $\uparrow\uparrow$   & $\uparrow\uparrow$   & $\uparrow\uparrow$ \\
+ UAD                  & $\uparrow$    & $\uparrow$    & $\uparrow$    & $\uparrow$ \\
+ CGG                  & $\uparrow$    & $\uparrow$    & $\uparrow$    & $\uparrow$ \\
+ RAFT + UAD           & $\uparrow\uparrow\uparrow$  & $\uparrow\uparrow\uparrow$  & $\uparrow\uparrow\uparrow$  & $\uparrow\uparrow\uparrow$ \\
+ RAFT + CGG           & $\uparrow\uparrow\uparrow$  & $\uparrow\uparrow\uparrow$  & $\uparrow\uparrow\uparrow$  & $\uparrow\uparrow\uparrow$ \\
\textbf{RAFT + UAD + CGG} & $\uparrow\uparrow\uparrow\uparrow$ & $\uparrow\uparrow\uparrow\uparrow$ & $\uparrow\uparrow\uparrow\uparrow$ & $\uparrow\uparrow\uparrow\uparrow$ \\
\bottomrule
\end{tabular}
\label{tab:main_results}
\end{table}

RAFT contributes the largest single reduction (it directly grounds factual claims in
retrieved evidence); UAD and CGG each add a smaller, complementary improvement; and the
three combine roughly additively, with the full system reducing hallucination
substantially over the baseline on all four benchmarks.

\subsection{Scale Dependence}

Applying the full RAFT + UAD + CGG system across the 4B, 8B, and 32B dense models, the
\emph{absolute} reduction in hallucination is roughly comparable across scales, even
though the larger models start from a lower baseline hallucination rate. This indicates
the interventions are largely orthogonal to base model capability: they add a fixed
factuality benefit on top of whatever the base model already achieves. The relative
mix shifts with scale — larger models benefit proportionally more from RAFT (they make
better use of retrieved evidence) and less from UAD (they are already better calibrated).

\subsection{Domain-Stratified Results}

Applying the full system, the largest \emph{relative} hallucination reductions occur in
the retrieval-amenable factual domains — medicine, law, and science — where RAFT's
grounding in retrieved evidence has the most to correct. Mathematics, which already had
the lowest baseline hallucination rate and depends less on external evidence, shows the
smallest relative reduction. The relative-reduction ordering therefore mirrors the
baseline-rate ordering: domains that hallucinate most at baseline improve most under
retrieval-augmented training.

\subsection{Citation Accuracy}

For the CGG citation head, both document-level and (the harder) span-level citation
accuracy improve with model scale: the 32B dense model identifies the correct source
document and span markedly more often than the smaller 4B and 8B variants. Span-level
accuracy is consistently lower than document-level accuracy, reflecting the greater
difficulty of localizing the exact supporting passage versus the correct document.

\subsection{Abstention Rate}

\begin{table}[H]
\centering
\caption{UAD abstention rate on answerable vs.\ unanswerable questions.}
\begin{tabular}{lrr}
\toprule
\textbf{Category} & \textbf{Abstention rate (\%)} & \textbf{Target} \\
\midrule
Genuinely unanswerable    & 74.2 & $>$70\% \\
Answerable (false abs.)   & 4.1  & $<$5\% \\
\bottomrule
\end{tabular}
\label{tab:abstention}
\end{table}

\subsection{Perplexity Impact}

A key concern is that anti-hallucination interventions may degrade general language
modeling quality. Measuring perplexity on the Pile and a held-out multilingual corpus
for the Zen-32B model before and after applying RAFT + UAD + CGG, we find the increase
in perplexity is under 1\% on both corpora. The interventions therefore improve
factuality without a meaningful cost to general language-modeling fluency.

%% -----------------------------------------------------------------------
\section{Related Work}
\label{sec:related}
%% -----------------------------------------------------------------------

Hallucination in language models has been studied from multiple angles. Retrieval-augmented
generation \cite{lewis2020retrieval} grounds model outputs in retrieved evidence.
RLHF \cite{ouyang2022training} aligns model behavior with human preferences, including
factual accuracy. Self-consistency \cite{wang2022selfconsistency} aggregates multiple
samples to reduce hallucination via majority vote. Our work differs in targeting
hallucination at the training, decoding, and output-structure levels simultaneously.

%% -----------------------------------------------------------------------
\section{Conclusion}
\label{sec:conclusion}
%% -----------------------------------------------------------------------

We have presented a three-component system — RAFT, UAD, and CGG — that substantially
reduces hallucination in the Qwen3-based Zen models across all tested benchmarks and
domains. The system delivers large relative reductions in hallucination with less than
1\% perplexity degradation, demonstrating that factual accuracy and fluency are not
fundamentally in conflict. RAFT contributes the bulk of the improvement, with UAD and
CGG adding complementary gains; we recommend deploying all three components together for
production use cases where factual reliability is critical.

\begin{thebibliography}{9}
\bibitem{lin2021truthfulqa}
S. Lin, J. Hilton, O. Evans.
\textit{TruthfulQA: Measuring How Models Mimic Human Falsehoods}.
ACL, 2022.

\bibitem{min2023factscore}
S. Min, K. Krishna, X. Lyu, et al.
\textit{FActScoring: Fine-grained Atomic Evaluation of Factual Precision in Long Form
Text Generation}. EMNLP, 2023.

\bibitem{li2023halueval}
J. Li, X. Cheng, W.X. Zhao, J.-Y. Nie, J.-R. Wen.
\textit{HaluEval: A Large-Scale Hallucination Evaluation Benchmark for Large Language
Models}. EMNLP, 2023.

\bibitem{lewis2020retrieval}
P. Lewis, E. Perez, A. Piktus, et al.
\textit{Retrieval-Augmented Generation for Knowledge-Intensive NLP Tasks}.
NeurIPS, 2020.

\bibitem{ouyang2022training}
L. Ouyang, J. Wu, X. Jiang, et al.
\textit{Training Language Models to Follow Instructions with Human Feedback}.
NeurIPS, 2022.

\bibitem{wang2022selfconsistency}
X. Wang, J. Wei, D. Schuurmans, et al.
\textit{Self-Consistency Improves Chain of Thought Reasoning in Language Models}.
ICLR, 2023.
\end{thebibliography}

\end{document}
